\clearpage\pagestyle{otMain}\setlength{\otBlocEspacement}{2.5mm}\fontsize{11}{15}\selectfont\corrigesChapitre{14}{Construire et comprendre les triangles}
\begin{blocLivre}{cor-14-1}\begin{corrige}{1}{Troisième angle}\demarche{Compléter la somme des trois angles.}
La somme vaut $180^\circ$ : $180-48-67=\res{65^\circ}$.
\label{fig-cor-14-1}\figureGeo[corrige]{tri-1}
\end{corrige}\end{blocLivre}
\begin{blocLivre}{cor-14-2}\begin{corrige}{2}{Un isocèle}\demarche{Partager le reste entre deux angles égaux.}
Les angles à la base sont égaux. Ils totalisent $180-44=136^\circ$, donc chacun vaut $136:2=\res{68^\circ}$.
\label{fig-cor-14-2}\figureGeo[corrige]{tri-2}
\end{corrige}\end{blocLivre}
\begin{blocLivre}{cor-14-3}\begin{corrige}{3}{Une construction à trois côtés}\demarche{Croiser deux arcs de compas.}
Tracer AB de 6 cm. Tracer un arc de centre A et de rayon 4 cm, puis un arc de centre B et de rayon 5 cm. Une intersection donne C ; relier. Il y a deux positions symétriques par rapport à (AB).
\label{fig-cor-14-3}\figureGeo[corrige]{tri-3}
\end{corrige}\end{blocLivre}
\begin{blocLivre}{cor-14-4}\begin{corrige}{4}{Aire directe}\demarche{Associer la bonne base et la bonne hauteur.}
$A=9\times4:2=\res{18\ \mathrm{cm}^2}$. Base et hauteur correspondent à un même sommet et à son côté opposé.
\label{fig-cor-14-4}\figureGeo[corrige]{tri-4}
\end{corrige}\end{blocLivre}
\begin{blocLivre}{cor-14-5}\begin{corrige}{5}{Reconnaître une médiane}\demarche{Utiliser le milieu indiqué.}
Elle joint le sommet A au milieu du côté opposé [BC]. C’est exactement la définition d’une $\res{\text{médiane}}$. Aucune perpendicularité n’est nécessaire.
\label{fig-cor-14-5}\figureGeo[corrige]{tri-5}
\end{corrige}\end{blocLivre}
\begin{blocLivre}{cor-14-6}\begin{corrige}{6}{Deux côtés et un angle}\demarche{Placer l’angle entre les côtés connus.}
Tracer AB. En A, tracer une demi-droite formant $60^\circ$ avec [AB). Y reporter 3 cm pour C, puis tracer [BC]. L’angle est construit au sommet A.
\label{fig-cor-14-6}\figureGeo[corrige]{tri-6}
\end{corrige}\end{blocLivre}
\begin{blocLivre}{cor-14-7}\begin{corrige}{7}{Un côté et deux angles}\demarche{Construire les angles aux bonnes extrémités.}
Tracer AB puis les demi-droites faisant 50 degrés en A et 60 degrés en B, dans le même demi-plan. Leur intersection est C. Le dernier angle vaut $180-50-60=\res{70^\circ}$.
\label{fig-cor-14-7}\figureGeo[corrige]{tri-7}
\end{corrige}\end{blocLivre}
\begin{blocLivre}{cor-14-8}\begin{corrige}{8}{Hauteur extérieure}\demarche{Chercher une perpendiculaire au support.}
Le pied H de la perpendiculaire se trouve au-delà de B. On prolonge la droite (AB), puis on trace la perpendiculaire à (AB) passant par C. Une hauteur porte sur la droite du côté, pas seulement sur le segment.
\label{fig-cor-14-8}\figureGeo[corrige]{tri-8}
\end{corrige}\end{blocLivre}
\begin{blocLivre}{cor-14-9}\begin{corrige}{9}{Cercle circonscrit}\demarche{Trouver un point équidistant des trois sommets.}
Tracer les médiatrices de [AB] et [AC]. Leur intersection O vérifie $OA=OB$ et $OA=OC$, donc les trois distances sont égales. Tracer le cercle de centre O passant par A ; il passe aussi par B et C.
\label{fig-cor-14-9}\figureGeo[corrige]{tri-9}
\end{corrige}\end{blocLivre}
\begin{blocLivre}{cor-14-10}\begin{corrige}{10}{Hauteur et aire}\demarche{Retrouver un facteur de la formule.}
$8h:2=24$, donc $4h=24$ et $h=\res{6\ \mathrm{cm}}$. On peut aussi doubler l’aire puis diviser par la base.
\label{fig-cor-14-10}\figureGeo[corrige]{tri-10}
\end{corrige}\end{blocLivre}
\begin{blocLivre}{cor-14-11}\begin{corrige}{11}{Aire et médiane}\demarche{Comparer bases et hauteurs.}
Les bases BM et MC ont la même longueur et les triangles ont la même hauteur issue de A. Leurs aires sont égales et totalisent 35 cm$^2$ : chacune vaut $\res{17{,}5\ \mathrm{cm}^2}$.
\label{fig-cor-14-11}\figureGeo[corrige]{tri-11}
\end{corrige}\end{blocLivre}
\begin{blocLivre}{cor-14-12}\begin{corrige}{12}{Triangle rectangle : hauteurs et cercle}\demarche{Repérer les perpendiculaires déjà présentes.}
(AB) et (AC), perpendiculaires, sont les hauteurs issues de B et de C ; l’orthocentre est $\res A$. Les médiatrices de [AB] et [AC] se coupent en O, $\res{\text{milieu de [BC]}}$ : c’est le centre du cercle circonscrit.
\label{fig-cor-14-12}\figureGeo[corrige]{tri-12}
\end{corrige}\end{blocLivre}
\begin{blocLivre}{cor-14-13}\begin{corrige}{13}{Un isocèle et ses droites}\demarche{Exploiter la symétrie d’un triangle particulier.}
C’est une médiane par définition de M. L’axe de symétrie du triangle isocèle passe par A et M ; il est perpendiculaire à (BC). C’est donc aussi la médiatrice de [BC] et la hauteur issue de A. Le centre du cercle circonscrit appartient à toutes les médiatrices, donc en particulier à \res{(AM)}.
\label{fig-cor-14-13}\figureGeo[corrige]{tri-13}
\end{corrige}\end{blocLivre}
\begin{blocLivre}{cor-14-14}\begin{corrige}{14}{Un tracé impossible}\demarche{Comparer le plus grand côté à la somme des deux autres.}
Pour 3, 4 et 8 : $3+4<8$, donc impossible. Pour 3, 4 et 7 : $3+4=7$, les points sont alignés, le triangle est aplati. Aucun des deux n’est un triangle non aplati.
\label{fig-cor-14-14}\figureGeo[corrige]{tri-14}
\end{corrige}\end{blocLivre}
\begin{blocLivre}{cor-14-15}\begin{corrige}{15}{Deux bases pour la même aire}\demarche{Calculer une aire commune avant de changer de base.}
L’aire vaut $12\times5:2=\res{30\ \mathrm{cm}^2}$. Avec la base BC, $10h:2=30$, donc $h=\res6$ cm. La médiane issue de C partage l’aire en deux : $\res{15\ \mathrm{cm}^2}$ chacune.
\label{fig-cor-14-15}\figureGeo[corrige]{tri-15}
\end{corrige}\end{blocLivre}
